\documentclass[a4paper]{article}
\usepackage{amsfonts}
\usepackage{amsmath}
\usepackage{amssymb}
\usepackage{graphicx}     

\newcommand{\Bgtan}{\operatorname{Bgtan}}

\begin{document}
  
\noindent \large Auteur: Siebe Mees \\
\noindent \large Studierichting: Informatica\\
\noindent \large Oefeningenreeks 2B nr. 6.6\\

\medskip

\normalsize


$2\Bgtan\dfrac{1}{2} - \Bgtan\dfrac{1}{7}$\\

We gebruiken de dubbele hoekformule: $\tan(2\alpha) = \dfrac{2\tan(\alpha)}{1 - \tan^2(\alpha)}$\\

We mogen deze formule toepassen omdat zowel $2\Bgtan\dfrac{1}{2}$ \text{ als } $\Bgtan\dfrac{1}{7}$ \ binnen het beeld van $\Bgtan(x)$ vallen, namelijk $\left(-\dfrac{\pi}{2}, \dfrac{\pi}{2}\right)$. \\

$ = \Bgtan\dfrac{2 \cdot \dfrac{1}{2}}{1 - \left(\dfrac{1}{2}\right)^2} - \Bgtan\dfrac{1}{7}$\\

$= \Bgtan\dfrac{1}{\dfrac{3}{4}} - \Bgtan\dfrac{1}{7}$\\

$= \Bgtan\dfrac{4}{3} - \Bgtan\dfrac{1}{7}$\\

$= \Bgtan\left(\dfrac{\dfrac{4}{3} - \dfrac{1}{7}}{1 + \dfrac{4}{3} \cdot \dfrac{1}{7}}\right)$\\

$= \Bgtan 1 = \dfrac{\pi}{4}$


\begin{thebibliography}{999}
%\bibitem{a} Referentie

\end{thebibliography}
\end{document} 